\section{背景与研究动机}

乳腺癌是全球女性发病率最高的恶性肿瘤。据世界卫生组织统计，2020年全球新增确诊病例超过226万例，占女性所有新发癌症病例的24.5\%。在中国，乳腺癌发病率同样呈持续上升趋势，已成为威胁女性健康的主要公共卫生问题。尽管诊断技术和治疗策略不断进步，乳腺癌的复发转移和耐药仍然是导致患者死亡的主要原因，而这一困境的根源在于肿瘤高度的分子异质性。

传统乳腺癌分类主要依赖组织病理学特征，如Bloom-Richardson分级、TNM分期等。这些指标提供了基础的风险评估框架，但无法充分解释为何组织学级别相似的患者在接受相同治疗后呈现截然不同的预后。临床实践表明，相同病理分期往往对应截然不同的分子机制和临床结局，这种现象促使研究者从分子层面寻找更精确的分类依据。

分子异质性的生物学基础源于肿瘤发生发展过程中的多层次调控异常。肿瘤细胞不仅在基因组层面存在拷贝数变异、点突变和结构重排，在表观遗传层面也表现出广泛的DNA甲基化异常和组蛋白修饰改变。这些调控层面的异常进一步影响转录组、蛋白组和代谢层的表达谱，最终形成具有不同分子特征的肿瘤亚群。准确识别这些亚群对于制定个体化治疗策略、改善患者预后具有重要指导意义。

基于基因表达谱的分子分型彻底改变了乳腺癌的临床管理范式。2009年，Parker等人首次提出PAM50分类系统~\cite{pam50}，通过检测50个核心基因的表达水平，将乳腺癌分为四个主要分子亚型：Luminal A、Luminal B、HER2-enriched和Basal-like。这一分类系统迅速被学术界和临床界广泛采用，成为乳腺癌分子分型的金标准。

Luminal A亚型表现为激素受体阳性、HER2阴性、Ki-67低表达，占全部乳腺癌的40\%至50\%，预后最优，10年生存率超过90\%，主要采用内分泌治疗。Luminal B亚型激素受体阳性但增殖指数高，或伴有HER2过表达，占15\%至20\%，预后介于Luminal A和HER2-enriched之间，通常需要化疗联合内分泌治疗。HER2-enriched亚型由HER2过表达驱动，激素受体阴性，占5\%至15\%，在抗HER2靶向治疗出现前预后较差，靶向治疗显著改善了该亚型患者的生存结局。Basal-like亚型为三阴性，增殖指数最高，占10\%至15\%，预后最差，主要依赖化疗，新型治疗策略如免疫检查点抑制剂和PARP抑制剂正在积极探索中。

从精准医学角度看，PAM50分子亚型分类直接指导患者的个体化治疗方案决策。Luminal亚型患者受益于内分泌治疗和CDK4/6抑制剂；HER2-enriched亚型需要抗HER2靶向治疗联合化疗；Basal-like亚型则以化疗为主。近年来，针对三阴性乳腺癌的免疫治疗和ADC药物也为该亚型带来了新的治疗希望。准确的分子亚型分类不仅具有重要的科学价值，更直接关乎患者的治疗效果和长期生存预期。

虽然基因表达数据在乳腺癌亚型分类中已建立坚实地位，但单一的转录组视角往往无法完整解释表型多样性。从生物学机制深度来看，肿瘤表型是多层次调控网络综合作用的结果。基因组层面，DNA序列变异、拷贝数变异、结构变异和胚系突变共同塑造肿瘤的遗传背景，BRCA1/2突变携带者具有更高的乳腺癌发病风险，且往往表现为Basal-like亚型。表观遗传层面，DNA甲基化、组蛋白修饰和非编码RNA调控形成复杂的表观遗传调控网络，启动子甲基化异常可导致抑癌基因沉默，是肿瘤发生的重要机制之一。转录调控层面，mRNA、lncRNA和miRNA的异常表达共同决定细胞的转录谱，转录因子网络的重编程影响细胞增殖、分化和凋亡等核心生物学过程。蛋白翻译层面，翻译后修饰如磷酸化、乙酰化、泛素化等和蛋白定位改变影响信号通路活性，蛋白质是大多数药物的直接作用靶点。代谢重编程层面，肿瘤细胞通过Warburg效应等代谢适应机制满足快速增殖的能量和物质需求，代谢酶的表达异常也是重要的预后标志物。肿瘤微环境层面，免疫细胞浸润、肿瘤相关成纤维细胞和血管生成等微环境因素显著影响治疗响应和患者预后。

仅使用RNA-seq数据进行分类存在固有局限。某些关键基因可能由于启动子甲基化而被转录沉默，其RNA表达量反映的是最终执行结果，而非上游调控因素。两个在RNA水平上表现相似的样本，其甲基化驱动的调控状态可能差异巨大，这将导致它们对靶向治疗或内分泌治疗的响应截然不同。整合甲基化数据可以帮助识别这些调控层面的差异，从而更准确预测治疗反应。RNA表达易受多种因素影响，技术层面包括测序深度差异、批次效应、文库制备质量和比对率；生物学层面包括细胞周期阶段、组织样本中肿瘤细胞比例、细胞类型组成和采样时间窗口。单一模态难以区分真正的表型差异与上述混杂因素导致的噪声。分子亚型并非离散的绝对分类，而是在转录谱空间中呈现连续分布。临床研究已观察到，相邻亚型样本在基因表达空间中相互重叠，形成边界模糊区。基于表达阈值的硬分类可能将边界样本错误归类，增加分类错误率。整合多组学信息可以提供额外的判别维度，有助于更准确识别边界样本的真实归属。RNA和甲基化数据并非完全独立，它们在生物学上存在调控与被调控的关系。启动子甲基化可以影响基因表达，而异常表达的转录因子也可以调控甲基化酶的表达。整合分析可以同时捕获基因表达变化的"是什么"和调控机制层面的"为什么"，从而实现对肿瘤表型的更全面刻画。

多组学融合的发展经历了三个主要阶段，每一阶段都对应不同的技术范式和应用场景。第一阶段是特征拼接，最直观的融合方式是将不同模态的特征向量直接拼接，形成超高维特征空间。这种方法实现简单、可解释性直接，是多组学融合的第一层对照，典型代表包括直接拼接RNA和甲基化矩阵后输入分类器。其优点在于不引入额外假设、适用范围广，特征空间中的交互信息可以被分类器自动学习，且实现简单、便于工程复现。局限性在于维度诅咒问题严重，在样本远小于特征的环境下易过拟合，无法显式建模模态间的交互关系，特征重要性解释困难。虽然如此，经过适当的特征选择后，拼接方法在乳腺癌分类中仍表现出竞争力。

第二阶段是表示学习，深度学习时代涌现了Autoencoder、VAE、MOFA等方法~\cite{mofa2018}，它们试图先学习每个模态的低维表示，再进行融合。MOFA是该阶段的代表性方法，通过概率因子模型建立可解释的潜因子结构，既能降低维度，又能对因子进行生物学解释。与直接拼接相比，表示学习的优势在于显式建模模态共享结构、降维压缩噪声、潜因子具有生物学可解释性。但在当前中小样本场景下，表示学习方法可能因模型自由度较高而不稳定。

第三阶段是结构化学习，近期方法倾向于显式引入生物网络结构，如蛋白质相互作用网络、基因调控网络和代谢通路网络等，通过图神经网络或路径分析获取更具有生物学意义的特征。图神经网络的优势在于能显式建模基因-基因、样本-样本或通路-通路之间的关系，因此在某些任务中能够捕获单纯矩阵模型难以表达的结构信息。但这类方法通常需要高质量的先验网络知识，且实现复杂度较高。

在本研究中，我们选择在前两阶段的框架内进行深入探索，通过对比拼接与MOFA两种策略的性能，评估表示学习在当前问题中的实际增益。这一选择既实用，又有助于揭示融合策略与任务特异性的关系，同时为后续扩展到结构化学习奠定方法论基础。

本研究的设计围绕三个核心科学问题展开，每个问题都针对多组学融合研究中的关键挑战。第一个问题是多组学融合是否能显著改进乳腺癌亚型分类的准确性与稳健性。我们的假设是通过整合表达与甲基化信息，模型能够从更多维度捕获亚型特异性的生物学特征，从而超越单一模态的性能基准。RNA数据提供了基因表达水平的直接视图，而甲基化数据则反映了上游表观遗传调控状态，两者的整合理论上应能同时捕获"基因表达结果"和"调控机制原因"，形成信息互补。我们预期Concat和Stacking方法因能利用跨模态互补信息而优于RNA-only。

第二个问题是不同融合策略在这一任务中的相对效能如何。我们预期不同融合策略各有优劣：早期融合实现简单但可能受到维度诅咒影响；潜表示融合可通过低维因子压缩噪声，但对模型参数和样本规模更敏感；晚期融合可通过分支级决策集成提升泛化能力，但同样对样本规模和超参数配置敏感。在当前中小样本条件下，我们预期简单方法可能因过拟合风险较低而表现更稳定。

第三个问题是特征筛选与预处理顺序在高维小样本分类中的关键程度。我们预期特征筛选与避免数据泄露的预处理顺序对分类性能的影响至关重要，其重要性甚至可能超过融合策略本身。在样本远小于特征的场景下，高方差特征筛选不仅能降低计算成本，更重要的是去除噪声特征、减少过拟合风险。同时，"先切分、后预处理"的严格流程是保证结果可信度的前提条件。

为了系统地回答这些问题，我们设计了主线对比实验、稳健性验证、误差解释和消融实验。主线对比实验在相同的数据分割与评估方案下，对RNA-only、Concat、MOFA、Stacking四种方案进行对标；稳健性验证采用重复分层交叉验证与统计检验评估结果可靠性；误差解释结合指标分布与置信区间分析性能差异来源；消融实验通过移除特定模态或特征选择步骤，分析各成分的边际贡献。本研究的技术路线遵循"数据加载-预处理-特征选择-模型训练-评估"的标准化流程，后续章节将围绕数据来源、特征处理、模型构建、实验设置和结果分析展开，力求在提供可复现框架的同时，通过严谨的方法设计回答上述科学问题。
